\section{\texorpdfstring{$\R^{n}$}{Rn}中的仿射子空间}

记号约定:$\left(\mathbf{e}_{i}\right)_{i=1}^{n}$是$\R^{n}$的典范基,$p\geq 1$.

\begin{proposition}{$p$维仿射子空间同胚于$p$维实数空间}{}
	对$\R^{n}$的每个$p$维非空仿射子空间$V$,存在$f\in\Aut_{\mathsf{Aff}}\left(\R^{n}\right)$满足如下条件:
	\begin{enumerate}
		\item $f\left(V\right)=\langle\mathbf{e}_{1},\ldots,\mathbf{e}_{p}\rangle$.
		\item $f\left(V\right)$是$\R^{n}$中同胚于$\R^{p}$的闭箱.
	\end{enumerate}
\end{proposition}

\begin{proposition}{$p$维仿射子空间是闭集}{}
	$\R^{n}$中每个非空的$p$维仿射子空间都是同胚于$\R^{p}$的闭集.
\end{proposition}

\begin{definition}{$\R^{n}$中的坐标轴}{}
	设$\left(D_{i}\right)_{i=1}^{n}$是$\R^{n}$中一族过$\mathbf{0}$的直线,对每个$1\leq i\leq n$有$\mathbf{e}_{i}\in D_{i}$.
	\begin{enumerate}
		\item 我们称$\left(D_{i}\right)_{i=1}^{n}$中的每一项是$\R^{n}$中的坐标轴.
		\item 当$n=2$时,分别称$D_{1}$和$D_{2}$为横轴和纵轴.
		\item 对每个$\mathbf{x}\in\R^{2}$和$i\in\left\{1,2\right\}$,$\mathbf{x}$的第$1$个坐标和第二个坐标分别称为$\mathbf{x}$的横坐标和纵坐标.
	\end{enumerate}
\end{definition}

\begin{proposition}{}{}
	设$\mathbf{a}\in\R^{n}$,$D$是$\R^{n}$中过点$\mathbf{a}$的直线,则存在$\mathbf{b}\in\R^{n}\setminus\left\{\mathbf{0}\right\}$使得$D=a+\R\mathbf{b}$,映射$f:\R\to D,t\mapsto\mathbf{a}+t\mathbf{b}$是同胚.
\end{proposition}

\begin{definition}{直线的方向向量和方向参数}{}
	设$\mathbf{a}\in\R^{n},\mathbf{b}\in\R^{n}\setminus\left\{\mathbf{0}\right\},D\coloneq\mathbf{a}+\R\mathbf{b}$.
	\begin{enumerate}
		\item 我们称$\mathbf{b}$是$D$的一个方向向量.
		\item 如果$\mathbf{b}=\sum\limits_{i=1}^{n}b_{i}\mathbf{e}_{i}$,则称$\left(b_{i}\right)_{i=1}^{n}$是$D$的方向参数.
	\end{enumerate}
\end{definition}

\begin{remark}{}{}
	如果$\mathbf{b}^{\prime}$也是$D$的方向向量,那么$\mathbf{b}\in\R^{\times}\mathbf{b}$.
\end{remark}

\begin{definition}{闭射线}{}
	设$\mathbf{a}\in\R^{n},\mathbf{b}\in\R^{n}\setminus\left\{\mathbf{0}\right\}$.
	\begin{enumerate}
		\item 我们称$\mathbf{a}+\R_{\geq 0}\mathbf{b}$是以$\mathbf{a}$为起点,$\mathbf{b}$为方向向量的(闭)射线.
		\item 如果$\mathbf{b}=\sum\limits_{i=1}^{n}b_{i}\mathbf{e}_{i}$,则称$\mathbf{a}+\R_{\geq 0}\mathbf{b}$是以$\mathbf{a}$为起点,$\left(b_{i}\right)_{i=1}^{n}$为方向参数的(闭)射线.
	\end{enumerate}
\end{definition}

\begin{proposition}{闭射线都连通,闭,同胚于$\R_{\geq 0}$}{}
	设$\mathbf{a}\in\R^{n},\mathbf{b}\in\R^{n}\setminus\left\{\mathbf{0}\right\}$,则$\mathbf{a}+\R_{\geq 0}\mathbf{b}$是同胚于$\R_{\geq 0}$的连通闭集.
\end{proposition}

\begin{definition}{方向相反的闭射线}{}
	设$\mathbf{a}\in\R^{n},\mathbf{b}\in\R^{n}\setminus\left\{\mathbf{0}\right\}$.我们称闭射线$\mathbf{a}+\R_{\geq 0}\mathbf{b}$和$\mathbf{a}+\R_{\geq 0}\left(-\mathbf{b}\right)$的方向相反.
\end{definition}

\begin{proposition}{$\R^{n}$中的直线能表示成两条方向相反的闭射线的并}{}
	设$\mathbf{a}\in\R^{n},\mathbf{b}\in\R^{n}\setminus\left\{\mathbf{0}\right\}$,则$\mathbf{a}+\R_{\geq 0}\mathbf{b}=\left(\mathbf{a}+\R_{\geq 0}\mathbf{b}\right)\cup\left(\mathbf{a}+\R_{\geq 0}\left(-\mathbf{b}\right)\right)$.
\end{proposition}

\begin{definition}{开射线}{}
	设$\mathbf{a}\in\R^{n},\mathbf{b}\in\R^{n}\setminus\left\{\mathbf{0}\right\}$.我们称$\mathbf{a}+\R_{>0}\mathbf{b}$是以$\mathbf{a}$为起点,$\mathbf{b}$为方向向量的开射线.
\end{definition}

\begin{proposition}{}{}
	设$\mathbf{a}\in\R^{n},\mathbf{b}\in\R^{n}\setminus\left\{\mathbf{0}\right\}$.
	\begin{enumerate}
		\item $\mathbf{a}+\R_{>0}\mathbf{b}$同胚于$\R_{>0}$,从而同胚于$\R$.
		\item 当$n\geq 2$时,$\mathbf{a}+\R_{>0}\mathbf{b}$不是$\R^{n}$的开集,是$\R^{n}$中包含$\mathbf{a}+\R_{>0}\mathbf{b}$的直线的闭集.
	\end{enumerate}
\end{proposition}

\begin{proposition}{过两点的直线的参数表示}{}
	设$\mathbf{x},\mathbf{y}\in\R^{n},\mathbf{x}\neq\mathbf{y}$,$D$是$\R^{n}$中过$\mathbf{x}$和$\mathbf{y}$的直线,那么$D=\left\{p\mathbf{x}+q\mathbf{y}\middle|\left(p,q\right)\in\R^{2},p+q=1\right\}$.
\end{proposition}

\begin{definition}{闭线段}{}
	设$\mathbf{x},\mathbf{y}\in\R^{n}$.我们称$\left\{u\mathbf{x}+v\mathbf{y}\middle|\left(u,v\right)\in\R_{\geq 0}^{2},u+v=1\right\}$是以$\mathbf{x},\mathbf{y}$为端点的(闭)线段.
\end{definition}

\begin{proposition}{}{}
	设$\mathbf{x},\mathbf{y}\in\R^{n}$.
	\begin{enumerate}
		\item $\left\{u\mathbf{x}+v\mathbf{y}\middle|\left(u,v\right)\in\R_{\geq 0}^{2},u+v=1\right\}$是$\R^{n}$中的连通紧集.
		\item 当$\mathbf{x}\neq\mathbf{y}$时,$\left\{u\mathbf{x}+v\mathbf{y}\middle|\left(u,v\right)\in\R_{\geq 0}^{2},u+v=1\right\}$同胚于$\left[0,1\right]$.
	\end{enumerate}
\end{proposition}

\begin{definition}{开线段,半开线段}{}
	设$\mathbf{x},\mathbf{y}\in\R^{n},\mathbf{x}\neq\mathbf{y}$.
	\begin{enumerate}
		\item 我们称$\left\{u\mathbf{x}+v\mathbf{y}\middle|\left(u,v\right)\in\R_{>0}^{2},u+v=1\right\}$是以$\mathbf{x},\mathbf{y}$为端点的开线段.
		\item 我们称$\left\{u\mathbf{x}+v\mathbf{y}\middle|\left(u,v\right)\in\R_{>0}^{2},u+v=1\right\}\cup\left\{\mathbf{y}\right\}$是在$\mathbf{x}$处开,在$\mathbf{y}$处闭的(半开)线段.
	\end{enumerate}
\end{definition}

\begin{proposition}{}{}
	设$\mathbf{x},\mathbf{y}\in\R^{n},\mathbf{x}\neq\mathbf{y},I=\left\{u\mathbf{x}+v\mathbf{y}\middle|\left(u,v\right)\in\R_{>0}^{2},u+v=1\right\}$.
	\begin{enumerate}
		\item $I$同胚于$\left(0,1\right)$,从而同胚于$\R$.
		\item $I\cup\left\{\mathbf{y}\right\}$同胚于$\left[0,1\right)$.
		\item $I\cup\left\{\mathbf{x}\right\}$同胚于$\left(0,1\right]$.
		\item $I,I\cup\left\{\mathbf{x}\right\},I\cup\left\{\mathbf{y}\right\}$都连通,它们的闭包都是$\left\{u\mathbf{x}+v\mathbf{y}\middle|\left(u,v\right)\in\R_{\geq 0}^{2},u+v=1\right\}$.
	\end{enumerate}
\end{proposition}

\begin{proposition}{仿射子空间的补集的拓扑性质}{}
	设$V$是$\R^{n}$的$p$维仿射子空间.
	\begin{enumerate}
		\item 如果$p<n$,那么$\R^{n}\setminus V$是$\R^{n}$的稠密开集.
		\item 如果$p<n-1$,那么$\R^{n}\setminus V$在$\R^{n}$中连通.
		\item 如果$p=n-1$,那么$\R^{n}\setminus V$在$\R^{n}$中有且仅有两个连通分支,二者都同胚于$\R^{n}$.
	\end{enumerate}
\end{proposition}

\begin{proposition}{}{}
	设$H$是$\R^{n}$中的仿射超平面,$\R^{n}\setminus H$的连通分支分别是$E_{1}$和$E_{2}$,那么$\closure\left(E_{1}\right)=E_{1}\cup H,\closure\left(E_{2}\right)=E_{2}\cup H$.
\end{proposition}

\begin{definition}{仿射超平面确定的开半空间和闭半空间}{}
	设$H$是$\R^{n}$中的仿射超平面,$\R^{n}\setminus H$的连通分支分别是$E_{1}$和$E_{2}$.
	\begin{enumerate}
		\item $E_{1},E_{2}$称为由$H$确定的开半空间.
		\item $E_{1}\cup H,E_{2}\cup H$称为由$H$确定的闭半空间.
	\end{enumerate}
\end{definition}